% !TeX root = ../main.tex
\begin{abstract}
We present a behavioral simulation framework for analog compute-in-memory inference, focusing on precision-retention trade-offs under realistic physical deployment constraints. The central result is a phase-change-memory (PCM) precision-retention frontier with a non-monotonic precision--accuracy curve: under the tested PCM UnitCell regime, 8-bit PCM achieves the highest fresh accuracy (77.60\%) with negligible drift, 4-bit PCM achieves comparable fresh accuracy (76.68\%) but violates a 1~pp deployment-drift SLA with a 4.01~pp one-day loss, and 6-bit PCM enters a device-to-device (D2D) sensitive transition zone with lower mean accuracy (68.44\%, seed-level std 6.03\%) but drift stability ($<$0.1~pp). To separate this physical frontier from algorithmic robustness, we first study an idealized AIHWKit IdealDevice regime. There, standard fixed-instance hardware-aware training reaches 87.28\% fresh-instance accuracy at 8-bit precision but collapses to 14.64\% at 4-bit; Ensemble Hardware-Aware Training (HAT), which resamples device-to-device mismatch masks during training, restores the 4-bit pure-quantization ablation to 86.16\%. These findings separate the algorithmic challenge of cross-instance robustness from the physical constraint that PCM deployment is gated by precision-dependent retention, with 8-bit emerging as the strongest overall operating point and 6-bit defining a quantization-D2D interaction regime rather than a Pareto optimum.
\end{abstract}
